%
% Chapter 1 — Introduction
%
\chapter{Introduction} \label{cap:intro}

\section{Context} \label{sec:intro_context}

In a media ecosystem increasingly marked by misinformation and the erosion of journalistic standards, \textit{Diário LX} emerges as a digital journalism initiative committed to editorial ethics, slow journalism, and proximity to local communities. The project is developed within the Journalism Trends Laboratory (\textit{Laboratório de Tendência no Jornalismo}), a structure affiliated with LIACOM (Applied Research Laboratory in Communication and Media) at the Escola Superior de Comunicação Social (ESCS) of the Instituto Politécnico de Lisboa (IPL).

The platform targets a young audience aged 18 to 35 years, with higher education backgrounds, and places a particular emphasis on the Greater Lisbon region. Editorially, it promotes visual journalism, long-form reporting, and investigative content, while deliberately avoiding clickbait-driven content strategies. The project has no commercial purpose and carries no advertising, operating instead with civic, educational, and scientific goals, co-funded by FCT (Fundação para a Ciência e Tecnologia).

\section{Motivation} \label{sec:intro_motivation}

A proof of concept was initially implemented using WordPress to validate the editorial structure and workflow. However, this approach revealed significant limitations in terms of architectural control, customization, and long-term scalability. The need to build a purpose-built platform that fully supports the editorial lifecycle — from content creation and review to publication and archival — motivated the development of this system from the ground up.

\section{Solution Overview} \label{sec:intro_solution}

\textit{Diário LX} is a web-based content management and publication platform
built to support the full editorial lifecycle of a digital journalism
outlet. The system is organized around two distinct areas:

\begin{itemize}
      \item \textbf{Backoffice}: A restricted editorial management interface, accessible only to authenticated users (Administrators, Editors, and Contributors), supporting content creation, revision, approval, and publishing workflows.
      \item \textbf{Public Website}: A reader-facing interface providing access to published articles, multimedia content (podcasts, videos, photo essays), organized by category, tag, and editorial section.
\end{itemize}

The system is built around an editorial workflow in which content progresses through
defined states — draft, pending review, published, and rejected — with each transition
guarded by the requesting user's role.

\section{Report Organization} \label{sec:intro_organization}

The remainder of this report is organized as follows:

\begin{itemize}
      \item Chapter~\ref{cap:background} presents the context, related work, and existing solutions analysed during the project.
      \item Chapter~\ref{cap:requirements} details the functional and non-functional requirements of the system.
      \item Chapter~\ref{cap:architecture} describes the system architecture and technology choices.
      \item Chapter~\ref{cap:datamodel} presents the data model.
      \item Chapter~\ref{cap:backend} covers the backend implementation.
      \item Chapter~\ref{cap:frontend} covers the frontend implementation.
      \item Chapter~\ref{cap:deployment} describes how the system is packaged and deployed to production.
      \item Chapter~\ref{cap:conclusion} concludes the report, summarising the delivered system and outlining directions for future work.
\end{itemize}
